\section{Background}

\subsection{Minimax Algorithm}

The Minimax algorithm is a foundational technique in adversarial game search. It models a two-player zero-sum game as a tree of alternating \textit{maximizing} and \textit{minimizing} nodes, where the maximizer (Player 1) seeks to maximize a heuristic evaluation score while the minimizer (Player 2) seeks to minimize it.

Formally, for a game state $s$ at depth $d$, the Minimax value is defined recursively as:

\begin{equation}
    \text{minimax}(s, d) =
    \begin{cases}
        h(s) & \text{if } d = 0 \text{ or } s \text{ is terminal} \\
        \max_{a \in A(s)} \text{minimax}(\text{result}(s,a),\, d{-}1) & \text{if maximizer's turn} \\
        \min_{a \in A(s)} \text{minimax}(\text{result}(s,a),\, d{-}1) & \text{if minimizer's turn}
    \end{cases}
\end{equation}

where $h(s)$ is a hand-crafted heuristic evaluation function, and $A(s)$ denotes the set of legal actions from state $s$.

\subsection{Alpha-Beta Pruning}

Alpha-beta pruning is an optimization of Minimax that eliminates branches provably unable to influence the final decision, without affecting the result. It maintains two bounds $\alpha$ (best guaranteed score for the maximizer) and $\beta$ (best guaranteed score for the minimizer), pruning any branch where $\alpha \geq \beta$.

In the best case, alpha-beta pruning reduces the effective branching factor from $b$ to $\sqrt{b}$, allowing the search to reach approximately twice the depth compared to plain Minimax within the same time budget. In practice, with good move ordering, the depth gain is substantial.

\subsection{Negamax Formulation}

For implementation simplicity, we adopt the \textbf{Negamax} formulation — a symmetric variant of Minimax where each node maximizes the negation of its children's scores:

\begin{equation}
    \text{negamax}(s, d, \alpha, \beta) =
    \begin{cases}
        h(s) & \text{if } d = 0 \text{ or } s \text{ is terminal} \\
        \max_{a \in A(s)} \bigl[-\,\text{negamax}(\text{result}(s,a),\, d{-}1,\, -\beta,\, -\alpha)\bigr] & \text{otherwise}
    \end{cases}
\end{equation}

This is mathematically equivalent to Minimax but requires only a single recursive case, which is both cleaner to implement and easier to augment with machine learning components.

\subsection{Hand-Crafted Heuristic Agent}

Our baseline agent uses a hand-crafted evaluation function $h(s)$ that combines multiple strategic features of the Othello board:

\begin{itemize}
    \item \textbf{Disc count}: weighted difference between the number of own and opponent discs.
    \item \textbf{Mobility}: number of legal moves available to each player.
    \item \textbf{Corner occupancy}: high-value corners that cannot be flipped once captured.
    \item \textbf{Stability}: discs guaranteed not to be flipped for the remainder of the game.
\end{itemize}

\subsection{Machine Learning for Game Playing}

Modern game-playing agents increasingly replace hand-crafted heuristics with learned functions. Neural networks trained on expert data can approximate the game's value function more accurately than manually designed formulas. A prominent example is AlphaGo \cite{silver2016mastering}, which combines Monte Carlo Tree Search (MCTS) with policy and value networks trained via self-play reinforcement learning.

In our work, we adapt this paradigm to Othello using a simpler, fully-connected (MLP) architecture trained via \textbf{supervised imitation learning} from a Minimax teacher.
